\chapter{UPPAAL model}\label{ch:uppaalmodel}
In order to decide upon a factory configuration in a given case, we must be able to evaluate configurations against each other. However, physically setting up configurations to compare them is in no way practical. Instead we choose to simulate individual factory configurations and virtually compare them. For simulation we need an abstract model of the modular factory and the items that are processed by it. With such a model in place, we should be able to simulate a specific set of items being produced on a specific configuration. In addition we would like to get the best schedule for producing a set of items on a configuration. A schedule is a sequence of actions, instructing how to go from having no items, to having produced all items. Comparing schedules between configurations producing the same items aids in picking, which configuration should be used for production. A configuration schedule may also be used as a jumping-off point for generating KLR code.

To model and simulate the modular factory system, we use the UPPAAL model-checker\cite{Larsen97uppaalin}. This is an integrated development environment used for modelling, simulating and verifying real-time systems. It allows us to model our system using timed automata in a combined graphical and programmatic manner. Once we have designed our model, we may use it to instantiate a configuration and a set of items. UPPAAL can then simulate the production of these items on the given configuration. In addition, the included UPPAAL model-checker allows us to query different properties about the setup. To us, the most interesting property is one of reachability; Is it possible to reach a state, where x amount of items have been produced? If this is the case, then UPPAAL may produce a fastest timed trace of transitions needed to reach this state. This evaluates to the temporaly shortest schedule to follow, in order to produce the given items on the given configuration.

UPPAAL has been used to model many different real-life systems, such as a lacquer\cite{so54514} and an industrial printer\cite{Igna2008}. From this research we take inspiration for our model. In these examples however, UPPAAL is used to model systems with only a single configuration. Our work differs in that we attempt to model a system of which there are many possible configurations, which need to be compared. Therefore we focus on keeping the module flexible and avoiding hard-coded values. Thus, many different configurations and items may be instantiated.

When modelling the real world, we need to take care to pick an appropriate abstraction level. Too high, and the model will not fit well with the real world. Too low and the model becomes very complex. The more complex a model is, the more likely it is that the search space will be larger. This severely increases the runtime when searching for a best schedule. Therefore, we naturally aim to hit a sweet spot between abstraction and complexity.

In this chapter we will explain how we modelled configurations and items in UPPAAL, so that we may simulate them. This will be a look at both the technical implementation, but also a discussion on when we have abstracted from the real world, and when we have not.  

\section{Recipe}\label{sec:recipe}
Looking into earlier projects developed using UPPAAL, an item being produced is often implemented as a sequence of actions, which must be taken in order to complete the item. Therefore items are often referred to as recipes. We will do the same from this point on. 

There are of course different ways of implementing a recipe. The simplest variant is a list of actions, which must be performed linearly. However, we find it more flexible to implement a recipe as an acyclic dependency graph. This is because we often do not care for the precise sequence of actions, just that some actions are performed before others. This is easily described using a dependency graph as seen in \cref{fig:dependency-graph}. This recipe describes that it needs to be hammered and screwed before it is packaged. Yet, it cares not in which order the hammering and screwing takes place. 

\begin{figure}[h]
\centering
\includegraphics[width=0.3\textwidth]{dependencygraph.pdf}
\caption{Dependency graph describing order of actions}
\label{fig:dependency-graph}
\end{figure}

We know that in some real life systems, a single item may be created through the aggregation of parts, which are individually produced. To escape this complexity, we only consider items, whose production is never distributed. That is, the production of an item starts at one single point and ends at another point never splitting up. As we represent recipes as dependency graphs, it also means that the same work can not be applied on a recipe several times, as this would produce cycles in the graph. 

In the following we describe how we handle recipes in our model through the use of the \emph{Recipe} and \emph{RecipeQueue} templates.


\subsection{Recipe}\label{subs:recipe}
It is possible to construct a functional dependency graph in UPPAAL using the graphical tools,  creating a new template for each new recipe type. However, as will be revealed in \cref{ch:configuration}, we wish to be able to instantiate a system and its recipes through python, instead of writing directly in UPPAAL. This requires that we edit the XML file in which our model is described. Having to generate a new template from scratch for each recipe type, would be very cumbersome. Instead we create a single recipe template that, depending on its instantiating parameters, functions according to a specific functional dependency graph. Thus we only have to generate code for instantiating templates.

Our recipe template can be seen in \cref{fig:recipe}. To describe the underlying functional dependency graph, the template is instantiated with an array of nodes. The \emph{Node} struct is shown in \cref{code:Node}. This reveals that each node knows what type of work it represents, the number of parents it has, its children's indexes in the node array and the number of children. Node A is parent of node B if B depends on A and the other way around for children. In addition a recipe is also instantiated with a unique id, the id of the module it begins at, as well as the direction it enters this module from. 

\lstinputlisting[language=C, caption=Node struct, captionpos = b, label={code:Node}, float]{codeRelated/UPPAAL/node.txt}

When a recipe begins processing, it is placed onto its start module, and we find the ids of nodes, which do not depend on any other node. The ids of these are placed into the local \emph{current\_nodes} array, and represent what works can be performed on the recipe. After this, the recipe is moved along the factory. Each time it meets a module, where it wishes to have work performed, it will handshake on its own private channel with the module to identify itself. Afterwards it will synchronize with the module again on one of the allowed \emph{work} channels. A work channel can be synchronized on, if one of the nodes in the \emph{current\_nodes} array, represent the work type corresponding to that work channel. This check is done by the \emph{is\_callable} function.

Before going back to the \emph{InProgress} location, a call to \emph{update\_current\_nodes} is made. The code of this function can be seen in \cref{code:updatecurrentnodes}. It will first collect all current nodes in the \emph{new\_nodes} array, except for \emph{called\_node}, as this is the node just worked on. It then runs over each of \emph{called\_node}'s children to decrement their \emph{number\_of\_parents} field by 1. If this field reaches 0, it means that the node no longer has any dependencies and can be worked on. Because of this it is added to the \emph{new\_nodes} array. Once all children have been updated, the contents of \emph{new\_nodes} are used to update \emph{current\_nodes}.  

\lstinputlisting[language=C, caption={The update\_current\_nodes function}, captionpos = b, label={code:updatecurrentnodes}, float]{codeRelated/UPPAAL/updatecurrentnodes.txt}

After the call to \emph{update\_current\_nodes} is finished, we call \emph{no\_more\_nodes}. This will set the local \emph{done} boolean to \emph{True}, if \emph{current\_nodes} nodes is empty. This means that the recipe can not be worked on any further and may be removed.

\subsection{RecipeQueue}\label{subs:recipequeue}

\begin{figure}[h]
\centering
\includegraphics[width=\textwidth]{recipe.pdf}
\caption{Recipe template}
\label{fig:recipe}
\end{figure}

Having all recipe instances compete for a place on the factory leads to a large state-space, when we ask UPPAAL to find the best timed trace. In order to get around this we enforce a certain order in which recipes may be placed into the factory.  This is done using a queuing system, which we implement using the \emph{RecipeQueue} template as seen in \cref{fig:recipequeue}.

The template is instantiated with an array of recipe id's indicating the order in which we wish the recipes produced. Once instantiated, the queue may have a recipe dequeued, which may then begin processing. This reduces the state space, as recipes are to begin in a specified order. We will not get into how to produce an efficient ordering of recipes here, as it will be brought up in \cref{ch:configuration}.

\begin{figure}[h]
\centering
\includegraphics[width=\textwidth]{recipequeue.pdf}
\caption{RecipeQueue template}
\label{fig:recipequeue}
\end{figure}
\section{Module}
\label{subs:module}
Since a configuration is made up of individual modules, it seems natural to model it as a set of synchronizing module instances. In this section we will first show off the initial design of the module template. Then discuss how we got more concrete with the final version.

\subsection{Early module template}
In our early design work we viewed the individual module in a very abstract way, inspired by initial observations of the FESTO CP-Learning Factory set-up at Aalborg University. It can receive and send recipes to any module connected to it. It may perform one type of work on a recipe and only one recipe may be located on it at any given time. Recipes may pass through modules without being worked on. 
These basic ideas turned into our first version of the module template, which can be seen in \cref{fig:earlymodule}.


\begin{figure}[h]
\centering
\includegraphics[width=\textwidth]{earlymodule.pdf}
\caption{Early version of the module template}
\label{fig:earlymodule}
\end{figure}

A module resides in the \emph{Idle} location when not processing a recipe. It leaves this location, when one of its up to four neighbouring modules passes on a recipe by synchronizing on the \emph{transport} channel. Once a recipe has been received, the module waits in the \emph{transporting} location for \emph{t\_time}. This simulates the time taken to transport the recipe across the module. Once the time has passed we may send the recipe to a neighbouring module. The guarding function \emph{is\_callable} makes sure, we only send to neighbours. This allows us to implement modules, which are only used for transporting recipes.

Instead of passing the recipe along, we may perform work on it by moving into the \emph{Processing} location. Here we wait for \emph{p\_time}, to simulate the time it takes to work the recipe. Before we can update the recipe to have its work performed, it must identify itself. This is done by synchronizing on the \emph{handshake} channel given by the value stored in the local \emph{lock} variable. The variable contains the unique id of the recipe, which most recently entered the module. It was received from the previous module over the global \emph{var} variable. This ensures that after a handshake, moving us to a committed location, the synchronizing on the \emph{work} channel will be with the correct recipe. Once work has been performed we may pass the recipe onto another module. 

Some modules may also allow for recipes to be removed from the system by transporting on the reserved \emph{transport[0]} channel. No actual module has the id 0. Instead an instance of the \emph{Remover} template is constantly calling on the \emph{transport[0]} channel. If synchronizing is established it does nothing with the received recipe. A module can only synchronize with the remover instance, if it has one of its neighbouring modules set to 0. 

\subsection{Final module template}
From further observations of the local FESTO set-up, we conclude that the above implementation is too abstract. In an actual factory, each module may perform several different types of work on a recipe. While only one recipe is worked at a time, several recipes may queue up on the module, waiting to be worked on or pass through. Some modules do not allow for recipes to pass through them, while they are working on another recipe. This feature is not enforced in the above implementation. In addition the time to transport a recipe across a module depend  on the dimensions of the module as well as the directions from which a recipe enters and leaves.

These features lead to some drastic changes in the module template. The biggest is that the template now gets split into three. \emph{ModuleQueue}, which controls the enqueuing and dequeuing of recipes. Then \emph{ModuleWorker} which performs work upon a recipe. Finally \emph{ModuleTransporter}, which transports recipes between modules. To get a complete module, we need an instance of each of these templates sharing the same module id. In the following we will explain each of these new templates.


\subsubsection{ModuleQueue}\label{subs:modulequeue}
One of the most important characteristics observed was the queuing nature of the system. Several recipes may be located on a module, but only one is being processed at a time. Some of the recipes queued may need to be processed, others just want to be passed along to a neighbouring module. Based on this behaviour we construct the \emph{ModuleQueue} template, which can be seen on \cref{fig:modulequeue}. An instance of this template has a fixed size of modules, which it may hold. Other modules may be able to use the \emph{enqueue} channel to move a recipe onto the module. In addition, recipes may be removed from it through a synchronization on its \emph{dequeue} channel.

As in the earlier version, a recipe is moved between instances by proxy of its id. The id is always stored in the local \emph{lock} variable.   

\begin{figure}[h]
\centering
\includegraphics[width=\textwidth]{modulequeue.pdf}
\caption{The ModuleQueue template}
\label{fig:modulequeue}
\end{figure}


\subsubsection{ModuleWorker}\label{subs:moduleworker}
The \emph{ModuleWorker} template can be seen in \cref{fig:moduleworker}. An instance of this template may apply several different types of work upon a recipe. In the idle location, the worker waits to synchronize with the module's \emph{ModuleQueue} instance. If the module can perform some work, it will move into the \emph{Working} location. Here it will wait for \emph{p\_time}, symbolising the time it takes to perform the work. When the wait is over, it will try to synchronize with the recipe through the \emph{handshake} and \emph{work} channels as in the earlier version. Arriving at the \emph{Done} \todo{Maybe call this location something other than Done. In general get names on all the locations} location, we may choose to work the recipe further if possible. Otherwise we perform a synchronization on the module's \emph{intern} channel. This passes the item over to the module's \emph{ModuleTransporter} instance. Synchronization on the intern channel may also occur, if the module cannot perform any work on the recipe. 

\begin{figure}[h]
\centering
\includegraphics[width=\textwidth]{moduleworker.pdf}
\caption{The ModuleWorker template}
\label{fig:moduleworker}
\end{figure}

\subsubsection{ModuleTransporter}
The \emph{ModuleTransporter} template can be seen in \cref{fig:moduletransporter}. Using an instance of this template, a module may transport an item onto a neighbouring module. If the item comes from the module's \emph{ModuleWorker} instance, then we move into the \emph{Selector} location by synchronizing on the module's \emph{intern} channel. We may however also move directly to the \emph{Selector} location by synchronizing with the module's \emph{ModuleQueue} instance over the \emph{dequeue} channel. This can however only be done, if the \cref{fig:moduletransporter} module allows us to pass a recipe through a module, while it is working on another recipe. This is indicated by the boolean \emph{pass\_through} variable, which is given at instantiation. To create a module only used for transportation, we set this variable to \emph{true} and do not create a \emph{ModuleWorker} instance for the module. 

Once in the \emph{Selector} location a recipe may go to \emph{Idle} or \emph{Transporting}. If the recipe is done, it has to go back to idle and be removed from the production line by synchronizing on the \emph{remove} channel, as opposed to the \emph{transport[0]} in the earlier version.

If the recipe is not done, it may look for a neighbouring module to be passed onto. A module may have up to four neighbours, one on each side. These are stored in the local \emph{next} array, their index indicating their placement relative to the module. When moving to \emph{Transporting} location, the recipe will choose a possible neighbour. In \emph{Transporting} we wait to simulate the time taken to transport a recipe over the module. This time varies according to where the recipe enters the module, and where it leaves. The different transport times are stored in the 4 by 4 multidimensional \emph{t\_time} array. Given that we at this point know the direction, which the recipe entered from \emph{origin} and the direction on which it will leave \emph{succ}, we can look up the exact time to wait in \emph{t\_time}. Once the wait is over we move to the \emph{Queuing} location. If the neighbour's queue is full, we wait here until there is room. When possible we synchronize with the neighbours instance of \emph{ModuleQueue} using the \emph{enqueue} channel. At the same time we use the local \emph{inverse} function to calculate, from which side the recipe enters the neighbouring module. This is sent with the global \emph{var2} variable and is later saved into the neighbour's \emph{origin} variable. 

\begin{figure}[h]
\centering
\includegraphics[width=\textwidth]{moduletransporter.pdf}
\caption{The ModuleTransporter template}
\label{fig:moduletransporter}
\end{figure}



%\begin{itemize}
%\item Modules are square or rectangle and have different physcial dimensions.
%\item Several types of work may be performed by the same module. \checkmark
%\item An item may pass through a module without have work performed on it. \checkmark
%\item Items queue up in a module and wait to be processed or just pass through. \checkmark
%\item Some modules allow for items to pass through while processing another item. Others do not. \checkmark
%\item A module may employ a seperate queue to store items that it is going to work on. 
%\item Some modules are only used for transport. \checkmark
%\item The time taken to transport an item over a module, depends on where the item enters and leaves. \checkmark
%\item An item may leave a module from any of its four sides, if there is another module ready to accept it. \checkmark
%\end{itemize} 

%On the list above, checkmarks indicate which characteristics we have chosen to implement directly in our model. While these observations were made from a specific system, we feel that they are general enough to apply to modular factories as a whole. 

%The physcial apsect of  modules are not enforced in our model. Two modules can be connected and pass items between each other, even though this may not physically be possible. Later in \cref{ch:configuration} we show, how we enforce some rules of physicality, when generating configurations outside of UPPAAL. We do not implement seperate module queues for items that need to be worked on. This would not have been difficult to implement. Yet, we see it as an edge case that would add onto our state space. Thus we choose to disregard it. The rest of the chracteristics we choose to implement in our model in one way or the other.


\section{Efficiency}\label{subs:efficiency}
\Cref{subs:recipequeue} illustrates how we make search for the shortest timed trace more efficient, by ordering how items are placed onto a factory. By applying such order, we are trying to keep down the global state space that needs to be searched, like how the bourgeois tries to keep down the proletariat. The lower the state space, the faster the search for the shortest timed trace.

This is not the only part of the model, where we combat the state space. We have strayed from going deeper into this until now, as it does not impact the overall functionality of the model. Yet it does greatly impact runtime. In the following subsections, we explain, how we have decreased our state space. 

\subsection{Urgencey}
The biggest positive impact on the state space has come through the use of committed locations and urgent channels. 

When a process enters a committed location, it ensures that globally no transition may be taken and no time may pass before the process leaves the location. This ensures that some processes stay atomic, such as handshake and work on item processes as described in \cref{subs:recipe}. This guards us from the interference of other processes, but also greatly reduces the state space, as we can not wait by making timed transitions in a commited location.

Urgency is an attribute that can be applied to different channels. Having a channel be urgent means that no time may pass, if we can transition from one process to another by synchronizing on the channel. In \emph{ModuleWorker} on \cref{subs:moduleworker} the \textit{Intern} channel is urgent. This means that, when the guards evaluate to true, then we transition by synchronizing on \emph{intern}. Again, not allowing the system the possibility of waiting reduces the state space.

There are transitions going between locations that we wish to make urgent, such as from \emph{Done} to \emph{Working} in the \textit{ModuleWorker} template, where there is not already some way of synchronizing on a channel. We solve this by creating a new simple template called \emph{Urgent}. A process instantiated from this will continuously try to synchronize on the urgent \emph{urg} channel. Thus we can add a synchronization on this channel to any non-urgent transition to make it urgent. 

A big difficulty with working with committed locations and urgent channels is that we sometimes need to be able to wait. This is the case when passing from \emph{ModuleWorker} to \emph{ModuleTransporter} through the intern channel as mentioned above. We can not safely make \emph{intern} urgent, if we may synchronize on it from \emph{ModuleWorker} but not \emph{ModuleTransporter}. This will lead to unwanted deadlock. Therefore, the guard on this particular transition states that it must only be taken, if the \emph{ModuleTransporter} is in idle and can communicate. We use this technique throughout the model, applying guards so that we may safely make our channels urgent. 

\subsection{Priority}
In order to shave more off the state space, we have also create priorities between channels using the built-in UPPAAL priority feature. This can be seen in \cref{code:cp}. By placing this order, we do not allow the model checker to synchronize on a certain channel, if a channel of a higher priority may be synchronized on. This helps guide us when there is an ambiguous choice of which channel to communicate on.

However, if we set up the same configuration in a version of the UPPAAL model without priorities and one with, then their initial parallel processes will not be bisimilar. This is because the model with priorities is restricted from transitions that the un-priotitized model may take. Yet, we find that this does not impact the fastest schedule found between the two different set ups. By design of the templates, none of the processes that make up the larger parallel system process will be able to make another transition, if they are waiting to communicate on a channel. Because of this, if two processes waiting to communicate are bypassed because of prioritization, they will stick around and wait. Resulting from the nature of the system, they will be able to communicate before the global clock is counted up, which only occurs as the result of transport and work on items. Because of this, our shortest time schedules end up having their final global clock values be equal, whether we prioritize or not.  

An issue is that even with channel priorities, we may have two transitions communicating on the same channel type, which again leaves us with an ambiguous order of execution. To get around this, we also prioritize system instances. If we ever get into the above situation, we pick the transition involving instances with the highest priorities. 

\lstinputlisting[language=C, caption=Channel priorities, captionpos = b, label={code:cp}, float]{codeRelated/UPPAAL/channelpriorities.txt}

\subsection{Initial Order}
A feature which UPPAAL templates lack is a user-defined constructor, which would allow us to execute code right when a process is made. We need this feature in templates such as \emph{ModuleQueue} in \cref{subs:modulequeue}, where we need to set up an array implementing the local queue. To get around this, we add an initial location to the template. To transition from this location, the instantiating functions have to be called.

The issue here is that with N processes having to be instantiated in this manner, there are N! ways of performing instantiation. This greatly increases our state space. To get around this, we use a queueing system similar to \emph{RecipeQueue} in \cref{subs:recipequeue}. This is implemented with the \emph{Initializer} template as seen in \cref{fig:initializer}. Each process, which needs to have code run upon instantiation, is given an instantiation id. When creating a procsses from the \emph{Initializer} template, we give it an array, which orders all these ids. The \emph{Initializer} process then starts to instantiate the instances in the order given to it by synchronizing on the \emph{initialize} channel. As the \emph{Initializer} process's main location is committed, nothing may occur until all needed instances have been synchronized with. Once all instantiations have been done, the \emph{Intanializer} is forced to move into a dead state, and we need not to worry about it anymore. From this point on the remaining  processes can transition as before. 

\begin{figure}[h]
\centering
\includegraphics[width=\textwidth]{init.pdf}
\caption{The \textit{Initialzier} template}
\label{fig:initializer}
\end{figure}


\section{Evaluation of Model}


rqasdad